% ECONOMICS_PROBLEM: {"id": "D63-653", "title": "Complexity of verifying epistemic EFX for additive goods", "jel_code": "D63", "source_id": "OP-0193"}
% !TeX program = xelatex
\documentclass[11pt,a4paper]{article}
\usepackage[margin=23mm]{geometry}
\usepackage{fontspec}
\setmainfont{Latin Modern Roman}
\usepackage{amsmath,amssymb,mathtools}
\usepackage{xcolor,enumitem,booktabs,longtable,array,xurl,hyperref}
\definecolor{muted}{HTML}{53636C}
\hypersetup{colorlinks=true,urlcolor=blue,linkcolor=blue}
\setlength{\parindent}{0pt}
\setlength{\parskip}{5pt}
\newcommand{\R}{\mathbb R}
\newcommand{\E}{\mathbb E}
\newcommand{\N}{\mathbb N}
\newcommand{\ind}{\mathbf 1}
\newcommand{\Var}{\operatorname{Var}}
\newcommand{\Cov}{\operatorname{Cov}}
\newcommand{\supp}{\operatorname{supp}}
\DeclareMathOperator*{\argmin}{arg\,min}
\DeclareMathOperator*{\argmax}{arg\,max}
\newcommand{\entrymeta}[3]{{\sffamily\small\color{muted}\textbf{\texttt{#1}}\quad|\quad #2\par #3\par}}
\newcommand{\Problemref}[1]{\texttt{#1}}
\newcommand{\JELref}[2]{\texttt{#2} (JEL \texttt{#1})}
\begin{document}
\section*{Complexity of verifying epistemic EFX for additive goods}
\textbf{Problem ID:} \texttt{D63-653}\quad \textbf{JEL:} \texttt{D63}\par
% BEGIN ECONOMICS STATEMENT
\label{problem:OP-0193}
% GAME_THEORY_PROBLEM: {"local_id": "GT-063", "original_number": 63, "kind": "P", "entry_kind": "statement_only", "published": false, "review_required": true, "category": "game_theory"}
\label{game:GT-063}
{\sffamily\small\color{muted}
\noindent\textbf{ID: GT-063.}\quad\textbf{Category:} Fair division complexity.
\quad\textbf{Kind: P.}\quad\textbf{Status:} Open in the cited source.
\par}
\subsection*{Definitions and mathematical statement}
Let \(N=\{1,\ldots,n\}\), \(n\geq1\), let \(M\) be a finite set of goods, and let
\(v_i(S)=\sum_{g\in S}a_{ig}\) for rational \(a_{ig}\geq 0\).
A complete allocation is an ordered partition of \(M\).
For this entry, agent \(i\) is EFX-satisfied by allocation \(Y\) if
\[
 v_i(Y_i)\geq v_i(Y_j\setminus\{g\})
 \quad
 \text{for every }j\ne i
 \text{ and every }g\in Y_j\text{ with }a_{ig}>0.
\]
An allocation \(A\) is epistemic EFX (\(\mathrm{EEFX}\)) if
\[
 \forall i\in N\ \exists\text{ complete allocation }Y^{(i)}:
 \quad Y^{(i)}_i=A_i
 \quad\text{and agent }i\text{ is EFX-satisfied by }Y^{(i)}.
\]

\paragraph{Decision problem.}
The input consists of \((n,M,(a_{ig}),A)\), explicitly encoded in binary,
where \(A\) is a complete allocation. Determine the computational
complexity of deciding whether \(A\) is \(\mathrm{EEFX}\).
In particular, is this decision problem polynomial-time solvable,
or is it NP-complete? It belongs to NP: the \(n\) allocations \(Y^{(i)}\)
are polynomial-size certificates whose defining inequalities can be
checked in polynomial time.

\subsection*{Short English statement}
How difficult is it to verify that a supplied allocation is EEFX?

\subsection*{Variants and scope boundaries}
Only positively valued goods are tested in the EFX condition here.
The valuation table is explicit and additive, rather than accessible
through a general set-function oracle. The goods problem is distinct
from the chore verification problem, for which [S1] proves hardness.

\subsection*{Source relationship and status}
[S1, Section 4, after Lemma 3] poses the goods verification question.
Its constructive procedure can provide certificates, which does not
give a verifier for arbitrary supplied allocations.

\subsection*{References}
\begingroup\footnotesize\raggedright
\noindent[S1] Ioannis Caragiannis, Jugal Garg, Nidhi Rathi,
Eklavya Sharma, and Giovanna Varricchio,
\emph{New Fairness Concepts for Allocating Indivisible Items},
IJCAI 2023, pp.~2554--2562; revised manuscript,
arXiv:2206.01710v3, 2025.
\url{https://arxiv.org/html/2206.01710v3}.
Locators: Definitions 1--2; Section 4, paragraph after Lemma 3;
Section 6, Theorem 7 (the separate chore hardness result).
\par\endgroup

\subsection*{Common conventions: Source mathematical conventions}

\textbf{Finite games.} For a finite set $A$, $\Delta(A)$ is its probability
simplex. Players choose independent mixed strategies in Nash equilibrium;
correlation is allowed only when explicitly part of the solution concept.
In a game with payoff functions $u_i$ and strategy profile $x$, an additive
$\varepsilon$ Nash equilibrium satisfies
\[
 u_i(x)\geq u_i(x_i',x_{-i})-\varepsilon
 \quad\text{for every player $i$ and every admissible deviation $x_i'$}.
\]
For numerical approximation, payoffs are normalized as locally specified.
An $\varepsilon$ payoff residual is distinct from distance at most
$\varepsilon$ to the set of exact equilibria. Point convergence, convergence
of empirical play, and convergence in distance to an equilibrium set are
also distinct.

\textbf{Computational representations.} Unless stated otherwise, a complexity
question takes rational data with binary encoding; polynomial time is measured
in total bit length. A fixed-accuracy polynomial algorithm is not necessarily
polynomial in $1/\varepsilon$, and neither is necessarily polynomial in
$\log(1/\varepsilon)$. Succinct games and value, demand, or sampling oracles
must declare their interfaces. Exact real solutions require an explicit
representation; an unspecified list of arbitrary real numbers is not a
finite computational output. A claimed algorithmic barrier is meaningful
only for its specified model and reductions.

\textbf{Dynamic games.} Histories, information, observation, and allowable
randomization are part of the model. A stationary strategy depends only on
the current state; a positional strategy in a deterministic graph game
chooses one action at each controlled vertex. Nash equilibrium at an initial
state and equilibrium after every history are different requirements.
An expectation of a pathwise $\liminf$ payoff, a $\liminf$ of expected averages,
a limiting expected average, and a uniform finite-horizon equilibrium are
different criteria. None is silently substituted for another.

\textbf{Allocation and mechanisms.} A complete allocation assigns every item
exactly once unless otherwise stated. Zero-valued goods, negative values,
capacity constraints, feasibility, lotteries, and copies of goods affect
fairness definitions and are specified locally. Dominant-strategy incentive
compatibility, Bayesian incentive compatibility, universal truthfulness,
and truthfulness in expectation impose different inequalities. Whenever
an objective ratio has zero denominator, its convention is stated or those
instances are excluded.

\textbf{Infinite-dimensional models.} Expectations, integrals, stochastic
equations, and optimization objectives require their local regularity,
integrability, filtrations, and admissible domains. A backward--forward
mean-field system is a boundary-value problem; it is not an initial-value
semiflow without an additional construction. Long-time, large-population,
and vanishing-noise limits require an order of limits and a convergence
topology. If a minimum need not be attained, the objective is its infimum
or supremum and the approximation requirement is stated separately.
% END ECONOMICS STATEMENT
\end{document}
